

\chapter{相关理论方法与数据集介绍}

\label{chap:related_theory}

\section{引言}

围绕交通资产清查这一应用背景，本章将聚焦于“目标检测增强”和“属性识别增强”两条技术主线：前者关注如何在复杂道路场景中更完整地发现资产目标，后者关注如何借助知识引导与数据补偿提升长尾异常属性的识别能力。本章将从目标检测基础模型、多模态大模型检测范式、GRPO 强化学习后训练方法、数据生成理论、监督微调技术以及实验所使用的数据集等方面展开说明，为本文后续研究的方法设计与实验分析提供理论与数据基础。

\section{目标检测相关理论基础}

\subsection{任务定义与评价指标}

设输入图像为 $I$，目标检测任务的真实标注集合记为

\begin{equation}
\mathcal{Y}=\{(b_i,c_i)\}_{i=1}^{N},
\end{equation}

其中 $b_i=(x_i,y_i,w_i,h_i)$ 表示第 $i$ 个目标的边界框，$c_i$ 表示对应的类别标签。检测器输出的预测集合可表示为

\begin{equation}
\hat{\mathcal{Y}}=\{(\hat b_j,\hat c_j,\hat s_j)\}_{j=1}^{M},
\end{equation}

其中 $\hat s_j$ 为预测置信度。边界框定位质量通常采用交并比（Intersection over Union, IoU）衡量，其定义为

\begin{equation}
\mathrm{IoU}(b,\hat b)=\frac{|b\cap \hat b|}{|b\cup \hat b|}.
\end{equation}

在此基础上，目标检测普遍采用 Precision、Recall、AP 以及 mAP 作为核心评价指标。其中，Precision 用于衡量预测结果的准确性，Recall 用于衡量真实目标的检出完整性，AP 用于刻画单类别在不同召回率水平下的综合性能，mAP 则是对所有类别 AP 的平均。对于交通资产清查任务而言，总体 mAP 虽然能够反映模型整体性能，但小目标召回率、密集目标检出率以及复杂区域中的漏检情况同样关键，因为这些指标直接关系到清查结果的完整性与工程可用性。

\subsection{目标检测常用方法}

从检测范式演进来看，目标检测方法大致可以分为双阶段检测器、单阶段检测器、基于 Transformer 的端到端检测器、开放词汇检测器以及多模态大模型检测器等类型。本文后续涉及的 Faster R-CNN、YOLO、DINO、Grounding DINO、DEIM 与 RexOmni 分别代表了上述不同技术路线，因此在本节中统一进行介绍。

\noindent\textbf{1）Faster R-CNN：典型双阶段检测器。}

Faster R-CNN 是双阶段目标检测方法的代表模型，其基本思想是先生成候选区域，再对候选区域进行分类与位置回归\cite{ren2015faster}。在结构上，该方法首先利用主干网络提取共享特征图，然后借助区域建议网络（Region Proposal Network, RPN）在特征图上生成候选框，最后将候选区域映射为固定长度的 RoI 特征，并送入分类分支与回归分支得到最终检测结果。

其训练目标通常写为

\begin{equation}
\mathcal{L}_{\text{FRCNN}}=\mathcal{L}_{\text{RPN}}^{\text{cls}}+\lambda_1\mathcal{L}_{\text{RPN}}^{\text{reg}}+\mathcal{L}_{\text{ROI}}^{\text{cls}}+\lambda_2\mathcal{L}_{\text{ROI}}^{\text{reg}},
\end{equation}

其中，$\mathcal{L}_{\text{RPN}}^{\text{cls}}$ 与 $\mathcal{L}_{\text{ROI}}^{\text{cls}}$ 分别用于学习候选框与 RoI 的分类能力，$\mathcal{L}_{\text{RPN}}^{\text{reg}}$ 与 $\mathcal{L}_{\text{ROI}}^{\text{reg}}$ 用于约束边界框位置。Faster R-CNN 的优势在于候选区域筛选和精细分类回归相互分离，因此通常具有较好的定位精度和较强的复杂背景鲁棒性；但其推理流程相对较长，实时性一般不及单阶段检测器。在本文场景下，该方法更适合作为高精度对照检测器或局部精检型基线。

\noindent\textbf{2）YOLO 系列：典型单阶段检测器。}

YOLO 系列将目标检测统一建模为单次前向传播中的密集预测问题，通过直接输出边界框位置、目标存在性及类别概率，实现端到端实时检测\cite{redmon2016yolo,yolov10}。与双阶段检测方法相比，YOLO 不显式生成候选区域，而是在不同尺度特征图上同时完成定位与分类，因此在速度与部署效率方面具有明显优势。

从技术演进角度看，YOLO 系列经历了从基于网格的直接回归、锚框机制、多尺度特征融合，到解耦检测头、无锚框设计以及后处理端到端化的发展过程。其训练目标可概括为

\begin{equation}
\mathcal{L}_{\text{YOLO}}=\lambda_{\text{box}}\mathcal{L}_{\text{box}}+\lambda_{\text{obj}}\mathcal{L}_{\text{obj}}+\lambda_{\text{cls}}\mathcal{L}_{\text{cls}},
\end{equation}

其中，$\mathcal{L}_{\text{box}}$ 用于约束边界框位置与尺度，$\mathcal{L}_{\text{obj}}$ 用于学习目标存在性，$\mathcal{L}_{\text{cls}}$ 用于完成类别判别。以 YOLOv10 为代表的新一代方法进一步引入一致性双分配等策略，试图在保持实时性的同时减弱对传统 NMS 的依赖\cite{yolov10}。

\noindent\textbf{3）DINO：基于 Transformer 的端到端检测器。}

DINO 是 DETR 路线的重要代表模型，其在端到端集合预测框架的基础上，通过改进去噪训练与查询初始化机制，有效缓解了原始 DETR 收敛较慢的问题，并显著提升了检测精度\cite{zhang2022dino}。该方法延续了 DETR 中“一个查询对应一个目标”的集合预测思想，利用 Hungarian matching 在预测集合与真实集合之间建立一一匹配关系，从而避免锚框设计和复杂后处理的依赖。

其匹配过程可表示为

\begin{equation}
\pi^{*}=\arg\min_{\pi}\sum_{i}\mathcal{C}(y_i,\hat y_{\pi(i)}),
\end{equation}

其中 $\mathcal{C}(\cdot)$ 为类别代价、位置代价与重叠度代价共同构成的匹配代价函数。匹配完成后，检测损失通常写为

\begin{equation}
\mathcal{L}_{\text{DINO}}=\sum_{i}\Big(\mathcal{L}_{\text{cls}}+\lambda_{1}\|b_i-\hat b_{\pi^{*}(i)}\|_{1}+\lambda_{2}\mathcal{L}_{\text{GIoU}}\Big).
\end{equation}

与原始 DETR 相比，DINO 通过 contrastive denoising、mixed query selection 以及 look forward twice 等策略，增强了查询学习能力与训练稳定性，使 Transformer 检测器在精度、可训练性与泛化性之间取得了更好的平衡。对于复杂道路场景中的多尺度交通资产识别，DINO 所具备的全局关系建模能力具有重要参考价值。

\noindent\textbf{4）Grounding DINO：开放词汇目标检测。}

Grounding DINO 将 DINO 检测器与大规模图文接地预训练结合起来，使模型能够根据文本提示完成开放词汇检测\cite{liu2024grounding}。与传统封闭类别检测器不同，该方法不再局限于固定类别集合，而是将文本编码后的语义表示与图像特征在 Transformer 中进行深层交互，使检测框与文本词元建立语义对齐关系。因此，模型既可以接受标准类别名称，也可以接受更加细粒度的自然语言描述，例如“damaged traffic sign”或“manhole cover near curb”。

若记图像查询特征为 $q_i$，文本词元特征为 $t_j$，则二者的语义相关性可表示为

\begin{equation}
s_{ij}=\sigma(q_i^{\top}t_j),
\end{equation}

其中 $s_{ij}$ 表示第 $i$ 个视觉查询与第 $j$ 个文本词元的对齐强度。训练阶段，Grounding DINO 仍沿用 DETR/DINO 的 bipartite matching 范式，但将传统封闭集分类损失替换为面向文本词元对齐的分类损失。其检测损失可概括为

\begin{equation}
\mathcal{L}_{\text{GDINO}}=\mathcal{L}_{\text{align}}+\lambda_{1}\mathcal{L}_{\text{bbox}}+\lambda_{2}\mathcal{L}_{\text{GIoU}},
\end{equation}

其中 $\mathcal{L}_{\text{align}}$ 表示基于查询--文本词元相似度的分类/对齐损失。依据原论文实现，模型先对查询与文本特征的点积 logits 计算 focal loss，再与边界框 L1 损失和 GIoU 损失共同用于匹配与最终优化。若记第 $i$ 个查询对应的文本监督为 $y_i^{\text{text}}$，则可进一步写为

\begin{equation}
\mathcal{L}_{\text{align}}=\sum_{i}\mathrm{FL}(s_i,y_i^{\text{text}}),\qquad s_i=\{s_{ij}\}_{j},
\end{equation}

其中 $\mathrm{FL}(\cdot)$ 为 focal loss。由此，Grounding DINO 在保持 DINO 端到端集合预测框架的同时，将“类别判别”扩展为“视觉查询与文本词元之间的跨模态对齐”。Grounding DINO 的主要优势在于具备开放词汇识别与零样本扩展能力，这对于交通资产清查中类别繁杂、标注难以穷尽的场景具有较高价值；但与此同时，其推理复杂度与文本提示敏感性也通常高于传统封闭集检测器。

\noindent\textbf{5）DEIM：面向快速收敛的改进匹配机制。}

DEIM 同样属于 DETR 路线上的改进方法，其目标是在保持实时检测性能的同时，缓解一对一匹配监督过于稀疏所带来的训练效率问题\cite{huang2025deim}。该方法提出 Dense O2O Matching，通过在标准数据增强框架下为单幅图像提供更多正样本监督，从而改善训练初期有效监督信号不足的问题；同时引入 Matchability-Aware Loss，对不同质量的匹配结果进行差异化约束，以降低密集匹配中低质量监督所带来的负面影响。

其中，DEIM 用于替代传统 VFL 的核心分类损失写为

\begin{equation}
\mathrm{MAL}(p,q,y)=
\begin{cases}
-q^{\gamma}\log(p)-(1-q^{\gamma})\log(1-p), & y=1,\\
-p^{\gamma}\log(1-p), & y=0,
\end{cases}
\end{equation}

其中 $p$ 表示分类置信度，$q$ 表示预测框与真实框之间的 IoU，$y\in\{0,1\}$ 表示前景/背景标签。与 Varifocal Loss 相比，MAL 用 $q^{\gamma}$ 显式刻画匹配质量，使低质量匹配在训练中仍能获得更充分的优化压力，从而提升 Dense O2O 监督的利用效率。结合边界框回归项后，其整体检测目标可写为

\begin{equation}
\mathcal{L}_{\text{DEIM}}=\mathcal{L}_{\text{MAL}}+\lambda_{1}\mathcal{L}_{\text{bbox}}+\lambda_{2}\mathcal{L}_{\text{GIoU}}.
\end{equation}

从原理上看，DEIM 的核心并不在于简单增加查询数量，而在于改进“哪些预测能够在训练中获得有效监督”的机制。也就是说，它一方面通过 Dense O2O 提高正样本密度，另一方面通过 MAL 抑制低质量匹配带来的噪声影响。对于需要在有限训练预算下快速得到较高性能检测器的工程场景，DEIM 具有较高的应用价值。本文后续在异构 Detection Agent 替换实验中引入该方法，正是为了验证主动感知框架对不同训练机制检测器的兼容能力。

\noindent\textbf{6）RexOmni：多模态大模型检测范式。}

RexOmni 将目标检测建模为多模态大模型上的自回归坐标生成任务\cite{rexomni,wang2024qwen2}。该模型建立在 Qwen2.5-VL-3B-Instruct 之上，通过复用视觉编码器、跨模态投影层和语言解码器，实现“图像输入 + 文本指令 $\rightarrow$ 坐标 token 序列输出”的统一建模。与传统检测器直接回归连续坐标不同，RexOmni 将归一化坐标量化到 $0\sim999$ 的离散空间中：

\begin{equation}
q(u)=\mathrm{round}(999u),\quad u\in[0,1],
\end{equation}

并通过

\begin{equation}
\hat u=\frac{q(u)}{999}
\end{equation}

将离散 token 近似映射回连续坐标。由此，一个边界框只需四个坐标 token 即可表示，从而降低了坐标学习难度并提升了空间表示效率。

RexOmni 的生成过程遵循标准自回归建模方式：

\begin{equation}
p_{\theta}(y|I,x)=\prod_{t=1}^{T}p_{\theta}(y_t|I,x,y_{<t}),
\end{equation}

其中 $I$ 为图像，$x$ 为文本提示，$y_t$ 为第 $t$ 个输出 token。该方法的重要特征在于：一是采用相对坐标 special token 表示空间位置；二是构建 grounding、referring、pointing 等多类数据引擎形成大规模视觉感知监督；三是采用“SFT 预训练 + GRPO 强化学习后训练”的两阶段训练范式提升几何定位精度与生成行为稳定性。对于本文研究而言，RexOmni 不仅是一个可直接使用的大模型检测器，更提供了一种将检测问题统一到生成式多模态建模框架中的思路。

在训练机制上，RexOmni 的一个关键环节是基于 Group Relative Policy Optimization（GRPO）的强化学习后训练\cite{shao2024deepseekmath,rexomni}。GRPO 可以视为 Proximal Policy Optimization（PPO）\cite{schulman2017ppo} 在大模型场景中的一种轻量化改造形式，其核心思想并非显式训练价值函数，而是对同一输入下采样得到的一组候选输出进行组内相对比较，利用标准化奖励构造优势项，从而在降低训练开销的同时提升输出排序与偏好优化能力。

给定输入 $(I,x)$，设当前策略 $\pi_{\theta}$ 采样得到 $G$ 个完整输出 $\{o_1,o_2,\ldots,o_G\}$，每个输出对应的标量奖励为 $r_i$。GRPO 首先构造组内相对优势：

\begin{equation}
A_i=\frac{r_i-\operatorname{mean}(r_1,\ldots,r_G)}{\operatorname{std}(r_1,\ldots,r_G)+\epsilon}.
\end{equation}

随后采用类似 PPO\cite{schulman2017ppo} 的截断目标进行优化：

\begin{equation}
\mathcal{J}_{\text{GRPO}}(\theta)=\frac{1}{G}\sum_{i=1}^{G}\min\Big(\rho_iA_i,\operatorname{clip}(\rho_i,1-\varepsilon,1+\varepsilon)A_i\Big)-\beta D_{\mathrm{KL}}(\pi_{\theta}\|\pi_{\mathrm{ref}}),
\end{equation}

其中 $\rho_i$ 表示新旧策略概率比，$\pi_{\mathrm{ref}}$ 通常取 SFT 阶段冻结得到的参考策略，$\beta$ 为 KL 正则项系数。该目标函数既保留了 PPO 类方法的稳定更新特征，又避免了训练 critic 所带来的额外复杂度。

对于 RexOmni 而言，引入 GRPO 的必要性主要体现在两个方面。其一，在 SFT 阶段，坐标 token 的训练主要依赖交叉熵损失，而离散 token 错误与真实几何偏差之间并不总是严格一致，容易出现“像素偏差很小但损失较大”或“仅有少量 token 错误但几何偏移严重”的问题。其二，在 teacher forcing 训练方式下，模型总是被动对齐真实输出序列，难以自主学习应当输出多少个检测框，因此推理阶段容易出现漏检、重复预测或框数量失控等行为问题。GRPO 恰好能够通过奖励设计同时约束几何质量与输出行为，使 RexOmni 在“坐标是否准确”和“输出行为是否合理”两个层面同时得到优化。

针对边界框任务，RexOmni 使用基于 IoU 的奖励聚合策略。设预测框集合为 $\hat{\mathcal{B}}=\{\hat b_i\}_{i=1}^{m}$，真实框集合为 $\mathcal{B}=\{b_j\}_{j=1}^{n}$，对每个真实框寻找与之 IoU 最大且类别一致的预测框，并将相应奖励记为 $r_j$。据此定义

\begin{equation}
\mathrm{Recall}=\frac{\sum_{j=1}^{n} r_j}{n},\qquad
\mathrm{Precision}=\frac{\sum_{j=1}^{n} r_j}{m},
\end{equation}

进一步构造 F1 风格的几何奖励：

\begin{equation}
r_{\text{box}}=\frac{2\cdot \mathrm{Precision}\cdot \mathrm{Recall}}{\mathrm{Precision}+\mathrm{Recall}+\epsilon}.
\end{equation}

这一奖励形式将边界框定位质量、类别正确性与重复预测惩罚统一在同一优化目标中，能够很好的完成目标检测的任务需求。

\section{数据生成与长尾样本增强理论基础}

\subsection{从经验风险最小化到分布补偿}

监督学习通常建立在经验风险最小化（Empirical Risk Minimization, ERM）基础上，即通过最小化训练集上的平均损失学习模型参数：

\begin{equation}
\hat R(\theta)=\frac{1}{N}\sum_{i=1}^{N}\ell(f_{\theta}(x_i),y_i).
\end{equation}

然而，在交通资产属性识别任务中，异常状态样本通常呈现显著长尾分布，高频正常样本占据训练集主体，而低频异常样本数量有限、类别分散。若直接依赖 ERM，模型往往更容易偏向头部类别，从而导致对关键异常属性的学习不足。

为缓解这一问题，数据增强与数据生成可被理解为对训练分布的主动补偿。Vicinal Risk Minimization（VRM）指出，可以围绕真实样本构造其邻域样本，从而扩展模型实际接触到的训练分布\cite{chapelle2001vrm}。其目标可写为

\begin{equation}
\hat R_{\text{vic}}(\theta)=\frac{1}{N}\sum_{i=1}^{N}\mathbb{E}_{(\tilde x,\tilde y)\sim \nu(\cdot|x_i,y_i)}\big[\ell(f_{\theta}(\tilde x),\tilde y)\big].
\end{equation}

Mixup 则给出了 VRM 的一种典型实现方式\cite{zhang2018mixup}：

\begin{equation}
\tilde x=\lambda x_i+(1-\lambda)x_j,\qquad
\tilde y=\lambda y_i+(1-\lambda)y_j,
\end{equation}

其中 $\lambda\sim\mathrm{Beta}(\alpha,\alpha)$。这一理论说明，合成样本的意义并不在于简单复制原始数据，而在于通过主动扩展样本分布、平滑类别边界来提升模型的泛化能力。

\subsection{生成式模型的数据合成原理}

当传统增强难以覆盖长尾异常的复杂变化时，需要显式学习数据分布并进行条件生成。一般地，数据生成问题可表示为在条件变量 $c$ 约束下，对真实分布进行近似建模：

\begin{equation}
\max_{\theta}\;\mathbb{E}_{(x,y)\sim p_{\text{data}}}\log p_{\theta}(x,y|c),
\end{equation}

其中 $c$ 可以由类别标签、文本描述、属性约束或领域知识共同构成。围绕这一目标，当前主流技术路线主要包括 GAN 与扩散模型两类。

GAN 通过生成器与判别器之间的对抗优化学习真实分布\cite{goodfellow2014gan}，其经典目标函数为

\begin{equation}
\min_{G}\max_{D}\;\mathbb{E}_{x\sim p_{\text{data}}}[\log D(x)]+\mathbb{E}_{z\sim p(z)}[\log(1-D(G(z,c)))].
\end{equation}

GAN 在生成速度和纹理细节方面具有一定优势，但训练稳定性和模式覆盖能力通常受限。扩散模型则通过逐步加噪和逐步去噪的方式逼近数据分布\cite{ho2020ddpm}，其前向扩散过程可写为

\begin{equation}
q(x_t|x_{t-1})=\mathcal{N}(\sqrt{1-\beta_t}x_{t-1},\beta_t\mathbf{I}),
\end{equation}

反向生成过程可写为

\begin{equation}
p_{\theta}(x_{t-1}|x_t,c)=\mathcal{N}(\mu_{\theta}(x_t,t,c),\Sigma_{\theta}(x_t,t,c)).
\end{equation}

由于扩散模型在多样性、条件控制能力和生成稳定性方面表现更优，当前高质量图像生成任务多采用扩散路线作为基础生成器。

\subsection{面向行业任务的数据生成约束}

对于交通资产属性识别任务而言，仅依赖通用生成模型往往难以保证样本真正可用。原因在于，异常属性样本不仅需要“视觉上合理”，还需要同时满足语义一致性、结构完整性以及字段约束正确性。因此，行业任务中的数据生成更适合采用“知识约束生成 + 自动质量筛选”的两阶段策略：首先基于领域知识构造条件变量 $c$，随后利用生成模型产生候选样本，并通过规则约束、识别模型或人工校核机制筛除低质量样本。

若记真实数据集为 $\mathcal{D}_{\text{real}}$，合成数据集为 $\mathcal{D}_{\text{syn}}$，则增强后的训练目标可表示为

\begin{equation}
\hat R_{\text{aug}}(\theta)=(1-\alpha)\hat R_{\text{real}}(\theta)+\alpha\hat R_{\text{syn}}(\theta),
\end{equation}

其中 $\alpha$ 用于控制合成数据对最终训练过程的影响程度。该表达式表明，合成数据的作用并非替代真实数据，而是对真实分布中稀缺但又关键的长尾区域进行有约束的补全。

\section{多模态大模型识别与 SFT 微调技术}

\subsection{多模态大模型的识别机理}

多模态大模型（Vision-Language Model, VLM）通常由视觉编码器、跨模态投影层以及大语言模型解码器构成\cite{liu2023visualinstruction,wang2024qwen2}。给定图像 $I$ 与文本指令 $x$，视觉编码器首先提取图像特征：

\begin{equation}
\mathbf{h}^{v}=E_{v}(I),
\end{equation}

随后通过投影层将视觉特征映射到语言模型可处理的统一表示空间：

\begin{equation}
\mathbf{z}=P(\mathbf{h}^{v}).
\end{equation}

最后，语言模型基于视觉表示 $\mathbf{z}$、文本指令 $x$ 以及历史输出 $y_{<t}$ 自回归地产生任务结果：

\begin{equation}
p_{\theta}(y|I,x)=\prod_{t=1}^{T}p_{\theta}(y_t|\mathbf{z},x,y_{<t}).
\end{equation}

这一统一生成接口使模型不仅能够进行开放式问答，还能够输出边界框、属性字段、异常描述等结构化结果。

\subsection{监督微调（SFT）的基本目标}

监督微调（Supervised Fine-Tuning, SFT）是当前大模型适配下游任务的基础方式，其本质是在高质量指令--响应数据上最大化目标输出序列的条件似然\cite{ouyang2022training,liu2023visualinstruction}。设训练样本记为 $(I,x,y^{*})$，则 SFT 的优化目标可表示为

\begin{equation}
\mathcal{L}_{\text{SFT}}=-\sum_{t=1}^{T}\log p_{\theta}(y_t^{*}|\mathbf{z},x,y_{<t}^{*}).
\end{equation}

对多模态任务而言，SFT 的作用至少体现在三个方面：第一，使模型学习任务要求的输出格式、字段顺序与表述规范；第二，使模型在特定领域术语、属性定义与标注规则上完成任务对齐；第三，为后续偏好优化或强化学习后训练阶段提供稳定的参考策略。LLaVA 等工作表明，在高质量视觉指令数据上进行 SFT，能够有效提升模型的多模态理解能力与任务迁移能力\cite{liu2023visualinstruction}。

\section{数据集描述}

本文实验基于 WHU-Infra3D 数据集开展。该数据集由移动测绘系统（Mobile Mapping System, MMS）采集，融合高分辨率全景图像与激光雷达点云，并提供多层次标注信息，用于支持城市道路基础设施的精细化感知与属性分析\cite{whu_infra3d_2025}。需要说明的是，尽管该数据集包含多模态数据及三维标注，本文仅使用其中的图像数据及对应的二维标注信息进行实验；同时，本文实际输入模型的并非原始全景图，而是由全景图经过四分投影切割后得到的 $4$ 张透视图像。

在实验设置上，本文选取上海（Shanghai）与武汉（Wuhan）两个子集构建跨域实验。上海数据采集于 2018 年，武汉数据采集于 2021 年，均由移动测绘系统获取，并包含分辨率为 $8192 \times 4096$ 的全景图像。为适配目标检测任务，原始全景图像首先被投影切割为 $4$ 张透视视图（rectilinear view），随后再进行标注组织、模型训练与测试。

在本文实际实验中，并不直接使用原始 ERP（Equirectangular Projection）全景图像作为检测输入，而是对 WHU-Infra3D 数据集中上海与武汉部分的全景图进行投影切割。具体而言，针对每幅宽高比为 $2{:}1$ 的全景图，先将其中像素坐标 $(u,v)$ 映射为单位球面上的方向角 $(\theta,\phi)$，再根据不同视角中心将球面方向重新投影到透视成像平面上，从而由 $1$ 张全景图生成 $4$ 张框幅透视图像。该处理能够有效缓解 ERP 图像在两极区域的几何拉伸问题，使道路资产在局部视图中的尺度、形状与常规相机成像更加一致，更适合后续目标检测模型学习。

若 ERP 图像尺寸为 $W \times H$，则像素到球面经纬角的映射可写为

\begin{equation}
\theta = 2\pi\left(\frac{u}{W}-\frac{1}{2}\right), \qquad
\phi = \pi\left(\frac{1}{2}-\frac{v}{H}\right).
\end{equation}

进一步地，球面方向可表示为单位向量

\begin{equation}
\mathbf{x}=
\begin{bmatrix}
\cos\phi\cos\theta\\
\cos\phi\sin\theta\\
\sin\phi
\end{bmatrix}.
\end{equation}

对于第 $i$ 个透视视图，记其外参变换为 $T_i$，则该方向在目标相机坐标系下可写为

\begin{equation}
\mathbf{X}_i=T_i^{-1}\mathbf{x},
\end{equation}

随后再通过透视投影模型 $\pi(\cdot)$ 得到图像平面坐标

\begin{equation}
\begin{bmatrix}
u_i\\
v_i
\end{bmatrix}
=
\pi(\mathbf{X}_i).
\end{equation}

如图~\ref{fig:erp_to_perspective} 所示，从 ERP 全景图到框幅图像的生成，本质上是“平面像素 $\rightarrow$ 球面方向 $\rightarrow$ 目标相机视角”的重投影过程。本文采用上述方式对数据进行统一预处理，最终得到武汉训练集 $6712$ 张透视图像，以及上海测试集 $4700$ 张透视图像；这两部分数据构成了本文实验实际使用的训练集与测试集。

\begin{figure}[htbp]

\centering

\begin{tikzpicture}[
    x=0.75cm,
    y=0.75cm,
    >=Stealth,
    font=\small,
    line width=0.9pt,
    every node/.style={inner sep=1.5pt}
]
    % ERP image
    \draw (0,0) rectangle (4.8,2.9);
    \fill[red] (2.9,1.6) circle (1.5pt);
    \node[above right] at (2.9,1.6) {$(u,v)$};
    \draw[->,gray] (2.2,1.25) -- (3.0,1.25);
    \node[gray,right] at (3.0,1.25) {$u$};
    \draw[->,gray] (2.35,1.55) -- (2.35,0.6);
    \node[gray,below] at (2.35,0.6) {$v$};
    \node[below=14pt] at (2.4,0) {\bfseries ERP 图像（2:1）};
    % Unit sphere
    \coordinate (C) at (9.6,1.45);
    \def\r{1.55}
    \draw (C) circle (\r);
    \draw[dashed,gray] ($(C)+(-\r,0)$) arc[start angle=180,end angle=360,x radius=\r,y radius=0.42];
    \draw[dashed,gray] ($(C)+(\r,0)$) arc[start angle=0,end angle=180,x radius=\r,y radius=0.42];
    \draw[dashed,gray] ($(C)+(0,\r)$) -- ($(C)+(0,-\r)$);
    \fill (C) circle (1.3pt);
    \node[below left] at (C) {$O$};
    \coordinate (P) at ($(C)+(50:1.45)$);
    \draw[->,blue,thick] (C) -- (P);
    \node[blue,right] at (P) {$\mathbf{x}(\theta,\phi)$};
    \draw ($(C)+(0.55,0)$) arc[start angle=0,end angle=50,radius=0.55];
    \node at ($(C)+(0.78,0.28)$) {$\phi$};
    \node[below=14pt] at ($(C)+(0,-\r)$) {\bfseries 单位球面};
    % Perspective image
    \draw[fill=blue!4] (14.1,0.25) rectangle (17.0,2.65);
    \fill[red] (15.95,1.7) circle (1.5pt);
    \node[right] at (15.95,1.7) {$(u_i,v_i)$};
    \draw[->,gray] (14.45,2.45) -- (15.1,2.45);
    \node[gray,right] at (15.1,2.45) {$u_i$};
    \draw[->,gray] (14.45,2.35) -- (14.45,1.75);
    \node[gray,above] at (14.45,2.35) {$v_i$};
    \node[below=14pt] at (15.55,0.25) {\bfseries 透视图像 $i$};
    % Arrows and annotations
    \draw[->] (5.35,1.45) -- (7.95,1.45);
    \node[above] at (6.65,1.55) {ERP 映射};
    \node[below] at (6.65,1.0) {$\theta=f(u)$};
    \node[below] at (6.65,0.55) {$\phi=f(v)$};
    \draw[->] (11.15,1.45) -- (13.95,1.45);
    \node[above] at (12.55,2.95) {外参变换};
    \node[above] at (12.55,1.45) {$\mathbf{X}_i=T_i^{-1}\mathbf{x}$};
    \node[below] at (12.55,1.0) {透视投影 $\pi$};
\end{tikzpicture}

\caption{ERP 全景图经球面重投影生成透视图像的原理示意图}

\label{fig:erp_to_perspective}

\end{figure}

在标注层面，数据集为每幅图像提供了面向 10 类道路基础设施的二维边界框标注，涵盖交通标志、路灯、信号灯、监控摄像头、路桩、消防栓、垃圾桶、井盖等典型城市设施。标注过程采用半自动生成与人工校验相结合的方式，以保证标注的一致性与准确性\cite{whu_infra3d_2025}。

尽管 WHU-Infra3D 同时提供点云语义标签、实例分割及属性状态等三维与高层语义信息，本文未使用相关模态，而是专注于二维视觉感知任务。这样可以在不引入几何模态辅助的前提下，评估模型仅依赖视觉输入时的目标发现与跨域泛化能力，也更贴近实际仅部署摄像头传感器的应用场景。

从数据分布来看，上海与武汉子数据集之间存在显著差异。上海数据主要采集于高密度城市区域，目标分布更为密集，遮挡现象较为严重；武汉数据则表现出更明显的光照变化与尺度变化。这种跨城市差异在目标外观和环境结构上形成了显著的域偏移\cite{whu_infra3d_2025}。

基于上述特点，本文采用“武汉训练、上海测试”的跨城市实验协议，以评估模型在分布迁移条件下的泛化能力。该设置能够有效模拟模型从已知城市迁移至新城市的实际应用场景，并避免在同分布数据上训练与测试带来的过拟合问题，从而使实验结果更具现实参考价值。

\section{本章小结}

本章围绕后续研究所需的理论基础与数据基础展开了系统梳理。首先，在目标检测相关理论基础部分，介绍了目标检测任务的基本定义与常用评价指标，并对 Faster R-CNN、YOLO、DINO、Grounding DINO、DEIM 与 RexOmni 等典型方法进行了统一归纳。其中，RexOmni 所体现的生成式检测范式以及其中融入的 GRPO 后训练思想，为本文后续多模态大模型检测研究提供了重要参考。其次，在数据生成与长尾样本增强理论部分，本章从分布补偿视角说明了长尾问题的形成原因，并进一步介绍了生成式模型的数据合成原理及面向行业任务的数据生成约束，为第四章的知识引导长尾样本构建奠定了理论基础。再次，在多模态大模型识别与 SFT 微调技术部分，本章说明了多模态模型的基本识别机理与监督微调目标，为后续属性识别增强方法提供了训练范式支撑。最后，本章对 WHU-Infra3D 数据集及其跨城市实验设置进行了说明，为第三章和第四章的实验验证提供了统一的数据基础。总体而言，本章内容构成了全文方法设计与实验分析的理论前提。
